\documentclass[12pt,a4paper]{report}

\usepackage[margin=1in]{geometry}
\usepackage[T1]{fontenc}
\usepackage[utf8]{inputenc}
\usepackage{lmodern}
\usepackage{microtype}
\usepackage{setspace}
\usepackage{array}
\usepackage{booktabs}
\usepackage{longtable}
\usepackage{tabularx}
\usepackage{xcolor}
\usepackage{enumitem}
\usepackage{listings}
\usepackage{fancyhdr}
\usepackage{titlesec}
\usepackage{caption}
\usepackage[hidelinks]{hyperref}

\definecolor{PhoneWatchBlue}{HTML}{1D4ED8}
\definecolor{PhoneWatchSlate}{HTML}{0F172A}
\definecolor{PhoneWatchGray}{HTML}{64748B}
\definecolor{CodeBackground}{HTML}{F8FAFC}

\hypersetup{
  pdftitle={PhoneWatch Project Report},
  pdfauthor={PhoneWatch Project},
  pdfsubject={Real-Time Phone Usage Detection},
  pdfkeywords={computer vision, YOLO, OpenCV, phone detection, Streamlit, Python}
}

\lstdefinestyle{phonewatch}{
  basicstyle=\ttfamily\small,
  backgroundcolor=\color{CodeBackground},
  frame=single,
  rulecolor=\color{PhoneWatchGray},
  breaklines=true,
  columns=fullflexible,
  keepspaces=true,
  showstringspaces=false
}

\lstset{style=phonewatch}

\setstretch{1.18}
\setlist[itemize]{topsep=4pt,itemsep=2pt,leftmargin=1.5em}
\setlist[enumerate]{topsep=4pt,itemsep=2pt,leftmargin=1.8em}

\pagestyle{fancy}
\fancyhf{}
\lhead{\textcolor{PhoneWatchGray}{PhoneWatch Project Report}}
\rhead{\textcolor{PhoneWatchGray}{\leftmark}}
\cfoot{\thepage}
\renewcommand{\headrulewidth}{0.4pt}

\titleformat{\chapter}
  {\normalfont\huge\bfseries\color{PhoneWatchSlate}}
  {\thechapter}{1em}{}
\titleformat{\section}
  {\normalfont\Large\bfseries\color{PhoneWatchSlate}}
  {\thesection}{0.8em}{}
\titleformat{\subsection}
  {\normalfont\large\bfseries\color{PhoneWatchSlate}}
  {\thesubsection}{0.8em}{}

\newcommand{\projectname}{PhoneWatch}
\newcommand{\projectrepo}{\url{https://github.com/emanamjad61/phonewatch}}
\newcommand{\codepath}[1]{\path{#1}}

\begin{document}

\begin{titlepage}
  \centering
  \vspace*{1.0cm}
  {\Huge\bfseries\color{PhoneWatchSlate} PhoneWatch\par}
  \vspace{0.35cm}
  {\Large\bfseries Real-Time Phone Usage Detection System\par}
  \vspace{1.2cm}
  {\large Project Report\par}
  \vfill
  \begin{tabular}{>{\bfseries}rl}
    Repository: & \projectrepo \\
    Primary Language: & Python \\
    Core Stack: & YOLO, OpenCV, MediaPipe, Streamlit \\
    Report Date: & April 22, 2026 \\
  \end{tabular}
  \vfill
  {\large Prepared for the PhoneWatch project repository\par}
  \vspace*{1.0cm}
\end{titlepage}

\pagenumbering{roman}

\chapter*{Abstract}
\addcontentsline{toc}{chapter}{Abstract}

\projectname{} is a computer-vision system designed to detect phone usage in real time from webcam, image, and video sources. The project combines object detection, contextual analysis, alert rendering, logging, training utilities, benchmarking, and a Streamlit dashboard into a single Python application. The detection pipeline uses a YOLO-compatible model to identify people and phones, then applies spatial context and hand-proximity analysis to decide whether a detected phone is likely being used. When usage is detected, the system renders visual overlays and records alert/session data for later review.

This report documents the problem context, objectives, requirements, architecture, implementation, dataset strategy, testing approach, operational workflow, limitations, and future development roadmap of the \projectname{} system. It is written as a formal project artifact for maintainers, evaluators, and future contributors.

\tableofcontents
\listoftables

\clearpage
\pagenumbering{arabic}

\chapter{Introduction}

\section{Project Overview}

\projectname{} is a Python-based computer vision project for monitoring visible phone usage in camera feeds. The system is intended for demonstration, research, and prototype monitoring scenarios where a user wants to detect when a phone appears to be held or used by a person. It provides:

\begin{itemize}
  \item real-time webcam inference;
  \item image and video file inference;
  \item contextual phone-use classification;
  \item visual alert overlays in meme and serious modes;
  \item session logging and alert history;
  \item dataset preparation and training utilities;
  \item benchmarking and diagnostics;
  \item a Streamlit dashboard for live monitoring and analytics.
\end{itemize}

The repository is organized around a single command-line entry point, \codepath{run.py}, with reusable modules under \codepath{src/}. Runtime outputs are written to \codepath{logs/}, trained models are stored under \codepath{models/checkpoints/}, and local datasets are stored under \codepath{data/}.

\section{Problem Statement}

Phone usage in monitored environments is difficult to detect reliably using simple object detection alone. A detector can identify a phone, but it cannot always determine whether the phone is being used, resting on a desk, visible in the background, or merely near a person. A useful system therefore needs both detection and context reasoning.

The central problem addressed by \projectname{} is:

\begin{quote}
How can a real-time computer-vision application identify phones and people, infer likely phone usage from visual context, and provide actionable alerts with usable logs and dashboard visibility?
\end{quote}

\section{Objectives}

The project objectives are:

\begin{enumerate}
  \item Detect phone and person objects in webcam, image, and video inputs.
  \item Infer whether a detected phone is likely in active use.
  \item Produce real-time annotated frames with bounding boxes and status overlays.
  \item Generate alert records when phone usage is detected.
  \item Support multiple alert presentation modes for formal and demonstration contexts.
  \item Provide a dashboard for live monitoring and historical analytics.
  \item Provide dataset, training, benchmarking, and test utilities for continued development.
  \item Maintain a clean, reproducible, and professional repository structure.
\end{enumerate}

\section{Scope}

The project is scoped as a prototype and applied computer-vision system. It is not a production surveillance product and should not be deployed in privacy-sensitive environments without additional consent, policy controls, and security review. Its primary purpose is to demonstrate an end-to-end detection workflow with practical engineering structure.

\chapter{Requirements}

\section{Functional Requirements}

\begin{longtable}{p{0.16\textwidth}p{0.76\textwidth}}
\caption{Functional requirements}\\
\toprule
\textbf{ID} & \textbf{Requirement} \\
\midrule
\endfirsthead
\toprule
\textbf{ID} & \textbf{Requirement} \\
\midrule
\endhead
FR-01 & The system shall process live webcam frames. \\
FR-02 & The system shall process image files and video files. \\
FR-03 & The system shall detect phone and person classes. \\
FR-04 & The system shall infer phone usage using contextual features, not only phone presence. \\
FR-05 & The system shall draw annotated output frames including bounding boxes and labels. \\
FR-06 & The system shall support meme and serious alert modes. \\
FR-07 & The system shall apply alert cooldowns to avoid excessive repeated alerts. \\
FR-08 & The system shall write session and alert logs under \codepath{logs/}. \\
FR-09 & The system shall expose a Streamlit dashboard for monitoring and analytics. \\
FR-10 & The system shall include dataset preparation and training utilities. \\
FR-11 & The system shall include smoke tests, unit tests, and benchmarking scripts. \\
\bottomrule
\end{longtable}

\section{Nonfunctional Requirements}

\begin{longtable}{p{0.16\textwidth}p{0.76\textwidth}}
\caption{Nonfunctional requirements}\\
\toprule
\textbf{ID} & \textbf{Requirement} \\
\midrule
\endfirsthead
\toprule
\textbf{ID} & \textbf{Requirement} \\
\midrule
\endhead
NFR-01 & The runtime workflow should be reproducible from documented commands. \\
NFR-02 & The codebase should separate source, datasets, model artifacts, and runtime logs. \\
NFR-03 & The system should avoid committing large generated datasets and output recordings. \\
NFR-04 & The CLI should provide clear user-facing error messages for missing files, models, and camera access. \\
NFR-05 & The dashboard should keep the detection worker isolated from the Streamlit rendering loop. \\
NFR-06 & The repository should be maintainable with a clear README, ignore policy, and test runner. \\
NFR-07 & The system should degrade gracefully when optional dashboard plotting dependencies are unavailable. \\
\bottomrule
\end{longtable}

\chapter{System Architecture}

\section{High-Level Architecture}

The \projectname{} application is organized as a layered pipeline. Inputs enter through the CLI or dashboard, frames are passed to the detection engine, usage context is inferred, alerts are generated, and outputs are logged or displayed.

\begin{table}[h!]
\centering
\caption{Architecture layers}
\begin{tabularx}{\textwidth}{p{0.20\textwidth}X}
\toprule
\textbf{Layer} & \textbf{Responsibility} \\
\midrule
Entry Points & \codepath{run.py}, \codepath{quick_test.py}, \codepath{benchmark.py}, and \codepath{dashboard/app.py}. \\
Detection Engine & \codepath{PhoneWatchEngine} loads the model, processes frames, draws annotations, logs sessions, and coordinates alert overlays. \\
Context Classifier & \codepath{PhoneUsageDetector} evaluates person-phone geometry and hand proximity. \\
Alert System & \codepath{AlertSystem} applies meme or serious overlays, cooldowns, sound alerts, and alert logging. \\
Data Pipeline & \codepath{DatasetManager}, \codepath{AugmentationPipeline}, and training utilities prepare data and train checkpoints. \\
Observability & Session CSVs, alert logs, dashboard charts, benchmark outputs, and test reports. \\
\bottomrule
\end{tabularx}
\end{table}

\section{Frame Processing Flow}

Each frame follows the same conceptual sequence:

\begin{enumerate}
  \item Capture or load a frame from webcam, video, or image source.
  \item Run YOLO inference for the target classes.
  \item Normalize detections into project-level dictionaries.
  \item Separate phone and person detections.
  \item Analyze spatial context between each phone and nearest person.
  \item Analyze hand proximity using MediaPipe hand landmarks when available.
  \item Produce one or more \codepath{UsageEvent} records.
  \item Draw object boxes, usage labels, alert overlays, and status HUD.
  \item Update session statistics and write session CSV rows.
  \item Emit cooldown-limited alert records when phone usage is active.
\end{enumerate}

\section{Repository Structure}

\begin{longtable}{p{0.26\textwidth}p{0.66\textwidth}}
\caption{Repository structure}\\
\toprule
\textbf{Path} & \textbf{Purpose} \\
\midrule
\endfirsthead
\toprule
\textbf{Path} & \textbf{Purpose} \\
\midrule
\endhead
\codepath{src/detect.py} & Main detection engine, model loading, inference, annotation, session logging. \\
\codepath{src/context.py} & Phone-use context classifier using bounding boxes and hand landmarks. \\
\codepath{src/alerts.py} & Alert overlay rendering, cooldown logic, and alert records. \\
\codepath{src/dataset.py} & Dataset acquisition, conversion, merge, split, and augmentation helpers. \\
\codepath{src/train.py} & YOLO training workflow and checkpoint management. \\
\codepath{src/demo_mode.py} & Presentation-oriented demo mode with overlays and recording. \\
\codepath{src/utils.py} & Shared configuration, logging, dataset analysis, diagnostics, and drawing helpers. \\
\codepath{dashboard/app.py} & Streamlit dashboard for live feed, controls, analytics, and model information. \\
\codepath{run.py} & Main CLI entry point for webcam, demo, image, video, train, setup, and dashboard. \\
\codepath{tests/} & Pytest suite covering dataset, context, alerts, and engine behavior. \\
\codepath{models/checkpoints/} & Trained and fallback model checkpoints. \\
\codepath{data/} & Local dataset workspace, intentionally ignored for large generated folders. \\
\codepath{logs/} & Runtime output directory, intentionally ignored except for \codepath{.gitkeep}. \\
\bottomrule
\end{longtable}

\chapter{Methodology}

\section{Object Detection}

\projectname{} uses a YOLO-compatible detector through the Ultralytics runtime interface. The model is configured to detect the project classes \codepath{phone} and \codepath{person}. The runtime parameters are loaded from \codepath{config.yaml}; the default confidence threshold is 0.5, the IoU threshold is 0.45, and the image size is 640 pixels.

The engine prefers the trained checkpoint \codepath{models/checkpoints/phonewatch_best.pt}. If that checkpoint is unavailable, it falls back to \codepath{models/checkpoints/yolov8n.pt} and then to the standard Ultralytics model resolution path.

\section{Contextual Phone-Use Classification}

Detecting a phone is not sufficient to decide whether the phone is being used. \projectname{} therefore performs context analysis through \codepath{PhoneUsageDetector}. The classifier combines:

\begin{itemize}
  \item nearest-person association;
  \item intersection-over-union and center-distance calculations;
  \item relative vertical position of the phone within the person box;
  \item heuristics for phones resting on surfaces;
  \item MediaPipe hand landmark proximity around the phone box.
\end{itemize}

This produces a structured \codepath{UsageEvent} with fields such as \codepath{in_use}, \codepath{confidence}, \codepath{method_used}, \codepath{person_id}, \codepath{phone_id}, and \codepath{reason}. The classifier is designed to reduce false positives from isolated phones and phones resting on desks while still flagging likely hand-held usage.

\section{Alert Strategy}

The alert system applies cooldown-limited alerts per person. When a usage event is active, \codepath{AlertSystem} renders either:

\begin{itemize}
  \item \textbf{meme mode}: a meme image overlay with a caption; or
  \item \textbf{serious mode}: a high-visibility red alert banner and bounding box emphasis.
\end{itemize}

Alert records are written to \codepath{logs/alert_log.csv}. The default alert mode is \codepath{meme}, and the default cooldown is 5 seconds.

\section{Dataset Strategy}

The dataset utilities support local and external YOLO-style data workflows. \codepath{DatasetManager} can prepare COCO phone-class examples, merge YOLO datasets, create train/validation/test splits, and write the \codepath{data.yaml} file expected by the training stack. \codepath{AugmentationPipeline} provides augmentation support for improving robustness.

Large generated data folders are intentionally excluded from git:

\begin{itemize}
  \item \codepath{data/raw/}
  \item \codepath{data/processed/}
  \item \codepath{data/augmented/}
\end{itemize}

This keeps the repository clean while preserving a reproducible data workspace.

\chapter{Implementation}

\section{Command-Line Interface}

The CLI in \codepath{run.py} is the main user entry point. It validates model paths, checks camera access, creates the detection engine, and dispatches to the selected mode.

\begin{longtable}{p{0.30\textwidth}p{0.58\textwidth}}
\caption{Primary commands}\\
\toprule
\textbf{Command} & \textbf{Purpose} \\
\midrule
\endfirsthead
\toprule
\textbf{Command} & \textbf{Purpose} \\
\midrule
\endhead
\codepath{python run.py webcam --camera 0} & Run live webcam detection. \\
\codepath{python run.py demo --camera 0} & Run presentation demo mode with enhanced overlays. \\
\codepath{python run.py video --input data/demo.mp4} & Run inference on a video file. \\
\codepath{python run.py image --input data/test.jpg} & Run inference on a single image. \\
\codepath{python run.py train --data data/processed/data.yaml} & Train using a YOLO data YAML file. \\
\codepath{python run.py dashboard} & Launch the Streamlit dashboard. \\
\codepath{python run.py setup} & Prepare datasets and first-time assets. \\
\bottomrule
\end{longtable}

\section{Detection Engine}

\codepath{PhoneWatchEngine} is responsible for:

\begin{itemize}
  \item loading configuration and model weights;
  \item resolving target class IDs;
  \item running inference per frame;
  \item parsing YOLO results;
  \item invoking the phone-use classifier;
  \item drawing live annotations;
  \item applying alert overlays;
  \item updating FPS and session statistics;
  \item writing frame-level session records.
\end{itemize}

The engine supports webcam, video, image, benchmark, and presentation-demo modes. It also provides backward compatibility through \codepath{PhoneWatchDetector}.

\section{Dashboard}

The Streamlit dashboard in \codepath{dashboard/app.py} provides an interactive interface for live detection and analytics. It uses a background worker thread for detection so that the Streamlit interface can render controls and charts independently of frame processing. Dashboard capabilities include:

\begin{itemize}
  \item start, stop, and reset controls;
  \item mode and confidence controls;
  \item live frame display;
  \item FPS, alert, and duration metrics;
  \item alert and usage analytics;
  \item model and dataset information panels;
  \item benchmark summaries.
\end{itemize}

\section{Training And Benchmarking}

Training is handled by \codepath{PhoneWatchTrainer}. It loads the project configuration, trains against a YOLO data YAML file, and produces checkpoints and training outputs. Benchmarking is handled by \codepath{benchmark.py}, which includes tests for inference speed, context classifier behavior, end-to-end latency, and memory profile.

\section{Configuration}

\codepath{config.yaml} centralizes project settings. Table~\ref{tab:config-summary} summarizes the most relevant values.

\begin{table}[h!]
\centering
\caption{Configuration summary}
\label{tab:config-summary}
\begin{tabularx}{\textwidth}{p{0.24\textwidth}X}
\toprule
\textbf{Category} & \textbf{Configured Values} \\
\midrule
Model & Weights path, confidence threshold 0.5, IoU threshold 0.45, image size 640. \\
Dataset & Raw, processed, augmented, train, validation, and label paths. \\
Training & 50 epochs, batch size 16, learning rate 0.001, automatic device selection. \\
Alerts & Meme mode by default, 5 second cooldown, meme image directory. \\
Dashboard & Streamlit port 8501. \\
Classes & \codepath{phone} and \codepath{person}. \\
\bottomrule
\end{tabularx}
\end{table}

\chapter{Testing And Quality Assurance}

\section{Test Coverage}

The pytest suite covers the most important internal contracts:

\begin{itemize}
  \item YOLO label conversion and dataset merge behavior;
  \item train/validation/test split integrity;
  \item context classification for desk, near-face, no-person, and multi-person cases;
  \item alert cooldown, mode toggle, serious overlay, and meme overlay behavior;
  \item engine initialization, single-frame processing, output format, and benchmark behavior.
\end{itemize}

The tests use synthetic frames and mocks where appropriate. This reduces dependence on a physical camera during automated testing.

\section{Smoke Testing}

\codepath{quick_test.py} provides an end-to-end smoke test without requiring a webcam. It either downloads a small sample video or generates a synthetic clip with OpenCV drawing, then runs the engine over a limited number of frames and saves an annotated sample frame.

\section{Manual Verification}

Manual verification should include:

\begin{enumerate}
  \item Running \codepath{python run.py --help} to verify CLI availability.
  \item Running \codepath{python quick_test.py} after dependency installation.
  \item Running \codepath{python run.py webcam --camera 0} with a real webcam.
  \item Running \codepath{python run.py dashboard} and validating dashboard controls.
  \item Reviewing \codepath{logs/alert_log.csv} and \codepath{logs/session_*.csv}.
\end{enumerate}

\section{Known Test Environment Constraint}

The cleaned repository intentionally excludes local virtual environments. A fresh checkout must run \codepath{./setup_env.sh} before full tests can be executed. Without the project dependencies, commands such as \codepath{pytest} and model inference will fail in a bare Python environment.

\chapter{Operational Workflow}

\section{Setup Workflow}

\begin{lstlisting}[language=bash,caption={Local setup workflow}]
chmod +x setup_env.sh run_tests.sh
./setup_env.sh
source .venv/bin/activate
\end{lstlisting}

\section{Live Detection Workflow}

\begin{lstlisting}[language=bash,caption={Live webcam detection}]
python run.py webcam --camera 0 --mode meme --confidence 0.5
\end{lstlisting}

\section{Dashboard Workflow}

\begin{lstlisting}[language=bash,caption={Launch dashboard}]
python run.py dashboard
\end{lstlisting}

\section{Demo Workflow}

\begin{lstlisting}[language=bash,caption={Presentation demo mode}]
python run.py demo --camera 0
\end{lstlisting}

\chapter{Privacy, Safety, And Ethical Considerations}

Computer-vision systems that observe people require careful handling. \projectname{} should be treated as a prototype and demonstration system unless additional governance is added. Responsible deployment requires:

\begin{itemize}
  \item clear consent from monitored individuals;
  \item visible notice that camera-based analysis is active;
  \item strict retention limits for logs and recordings;
  \item access controls for dashboards, logs, and model outputs;
  \item review of local laws and organizational policies;
  \item procedures for handling false positives and false negatives.
\end{itemize}

The system currently writes local logs and can record demo video output. Operators should avoid collecting or retaining unnecessary personal data.

\chapter{Limitations}

\section{Technical Limitations}

\begin{itemize}
  \item Detection accuracy depends on the model checkpoint and training data quality.
  \item Phone-use classification is heuristic and can produce false positives or false negatives.
  \item Hand-landmark analysis can be affected by occlusion, lighting, motion blur, and camera angle.
  \item Real-time performance depends on hardware, device selection, frame size, and model size.
  \item The dashboard is intended for local use and does not include production authentication.
  \item The repository excludes large datasets; reproducibility of training requires rebuilding or restoring local data.
\end{itemize}

\section{Operational Limitations}

\begin{itemize}
  \item Webcam access requires operating-system camera permission.
  \item Only one process should use the same camera at a time.
  \item The default alert cooldown may need tuning for different environments.
  \item Logs are local runtime outputs and are not automatically rotated or encrypted.
\end{itemize}

\chapter{Future Work}

Recommended future improvements include:

\begin{enumerate}
  \item Add a reviewed license file.
  \item Add CI after enabling GitHub token workflow permissions.
  \item Add structured model cards for checkpoints.
  \item Add a dataset versioning strategy using a tool such as DVC or external storage.
  \item Add configurable camera source selection in the dashboard.
  \item Add role-based access control if the dashboard becomes network accessible.
  \item Add alert export formats such as JSON Lines and Parquet.
  \item Add calibration tools for confidence thresholds and cooldown settings.
  \item Add a small reproducible evaluation subset with documented metrics.
  \item Add privacy controls for redaction or non-retention of raw frames.
\end{enumerate}

\chapter{Conclusion}

\projectname{} provides a complete prototype for real-time phone-use detection. It combines model inference, context-aware classification, alert rendering, logging, dashboard monitoring, training, benchmarking, and testing into a maintainable Python codebase. The architecture is modular enough for continued improvement, while the cleaned repository structure supports professional review and future collaboration.

The current implementation is best understood as a strong applied computer-vision project foundation. With further work on evaluation datasets, privacy controls, production hardening, and CI, it can evolve from a demonstration system into a more reliable operational tool.

\appendix

\chapter{Appendix A: Command Reference}

\begin{longtable}{p{0.42\textwidth}p{0.48\textwidth}}
\caption{Command reference}\\
\toprule
\textbf{Command} & \textbf{Description} \\
\midrule
\endfirsthead
\toprule
\textbf{Command} & \textbf{Description} \\
\midrule
\endhead
\codepath{./setup_env.sh} & Create local virtual environment and install dependencies. \\
\codepath{source .venv/bin/activate} & Activate local environment. \\
\codepath{python run.py --help} & Show CLI usage. \\
\codepath{python run.py webcam --camera 0} & Run real-time webcam mode. \\
\codepath{python run.py demo --camera 0} & Run presentation demo mode. \\
\codepath{python run.py video --input <file>} & Run video-file inference. \\
\codepath{python run.py image --input <file>} & Run image inference. \\
\codepath{python run.py train --data data/processed/data.yaml} & Start training. \\
\codepath{python run.py dashboard} & Launch Streamlit dashboard. \\
\codepath{python quick_test.py} & Run webcam-free smoke test. \\
\codepath{python benchmark.py} & Run benchmark suite. \\
\codepath{./run_tests.sh} & Run pytest suite with coverage report. \\
\bottomrule
\end{longtable}

\chapter{Appendix B: Key Runtime Artifacts}

\begin{longtable}{p{0.36\textwidth}p{0.54\textwidth}}
\caption{Runtime artifacts}\\
\toprule
\textbf{Artifact} & \textbf{Purpose} \\
\midrule
\endfirsthead
\toprule
\textbf{Artifact} & \textbf{Purpose} \\
\midrule
\endhead
\codepath{logs/alert_log.csv} & Cooldown-limited phone-use alert records. \\
\codepath{logs/session_*.csv} & Per-frame session summaries and usage events. \\
\codepath{logs/screenshots/} & Optional screenshots saved during detection. \\
\codepath{logs/demo_recording_*.mp4} & Demo mode recordings. \\
\codepath{logs/coverage_report/} & HTML coverage output from tests. \\
\codepath{models/checkpoints/phonewatch_best.pt} & Preferred trained PhoneWatch checkpoint. \\
\codepath{models/checkpoints/phonewatch_best.onnx} & ONNX export of the trained checkpoint. \\
\codepath{models/checkpoints/yolov8n.pt} & Fallback YOLO checkpoint. \\
\bottomrule
\end{longtable}

\chapter{Appendix C: Repository Publication}

The project has been published as a private GitHub repository at:

\begin{center}
\projectrepo
\end{center}

The repository metadata describes the project as real-time phone-use detection with YOLO, OpenCV, contextual alerts, and a Streamlit dashboard. The configured topics are Python, computer vision, YOLO, OpenCV, Streamlit, and phone detection.

\end{document}
